\paragraph{Ridge Regression}

\subparagraph{Hạng tử điều chỉnh}
Công thức hạng tử điều chỉnh trong Ridge Regression có dạng:
\begin{equation}
  \frac{1}{2}\|\mathbf{w}\|_2
\end{equation}
Hạng tử này chính là tổng bình phương các trọng số, còn được gọi là chuẩn
$L_2$ của vector trọng số.
Việc đưa hạng tử này vào hàm lỗi có tác dụng làm "phạt" các giá trị trọng số
lớn.

\subparagraph{Hàm mục tiêu}
Khi đó, hàm mục tiêu của Ridge Regression là:
\begin{equation}
  E(\mathbf{w}) = \frac{1}{2} \|\mathbf{Xw} - \mathbf{t}\|^2 + \frac{\lambda}
  {2}\|\mathbf{w}\|_2.
\end{equation}
\subparagraph{Nghiệm giải}
Ta giải được nghiệm của Ridge Regression là:
\begin{equation}
  \mathbf{w} = (\mathbf{X}^{\intercal}\mathbf{X} + \lambda \mathbf{I})^{-1}
  \mathbf{X}^{\intercal}\mathbf{t}.
\end{equation}
So với nghiệm của hồi quy tuyến tính thông thường:
\begin{equation}
  \mathbf{w} = (\mathbf{X}^{\intercal}\mathbf{X})^{-1} \mathbf{X}^{\intercal}
  \mathbf{t}.
\end{equation}
Ridge Regression bổ sung thêm $\lambda \mathbf{I}$ vào ma trận $\mathbf{X}^
{\intercal}\mathbf{X}$.
Điều này mang lại hai lợi ích quan trọng:
\begin{itemize}
  \item \textbf{Tránh ma trận suy biến}: Khi $\mathbf{X}^{\intercal}\mathbf{X}
  $ không khả nghịch, việc cộng $\lambda \mathbf{I}$ giúp ma trận luôn khả
  nghịch.
  \item \textbf{Ổn định nghiệm}: Giảm độ nhạy của nghiệm đối với nhiễu trong
  dữ liệu.
\end{itemize}

\subparagraph{Diễn giải trực quan}
Có thể hiểu Ridge Regression là bài toán tối ưu với ràng buộc:
\begin{equation}
  \min_{\mathbf{w}} \|\mathbf{Xw} - \mathbf{t}\|^2 \quad \text{với} \quad \|
  \mathbf{w}\|^2 \leq C.
\end{equation}
Tức là ta tìm nghiệm vừa khớp dữ liệu tốt, vừa nằm trong một "quả cầu" giới
hạn độ lớn của trọng số. Điều này giúp mô hình tránh học các hệ số quá lớn gây
overfitting.
